\begin{frame}{Tentative results}{MRT-calibrated simulator (\texttt{positive/negative} variant): RL vs.\ randomization}

\begin{columns}[T,onlytextwidth]
\column{.46\textwidth}
\centering
\includegraphics[width=\linewidth]{results_cae.pdf}

\column{.50\textwidth}
\scriptsize
\begin{tabular}{@{}lccc@{}}
\toprule
Policy & Mean $Y$ & \% sent & $\widehat e$ at wk 36 \\
\midrule
Micro-query + RS ($\bar\gamma{=}0.5$) & 5.444 & 53.1\% & --$^*$ \\
Randomized $p=0.5$                    & 5.406 & 50.0\% & 1.13 \\
\bottomrule
\end{tabular}

\vspace{.4em}
\normalsize
\begin{block}{Also tracked}
Belief-state calibration against simulated truth; sensitivity to a misspecified DAG.
\end{block}
\end{columns}

\vspace{.4em}
\begin{alertblock}{What would count as evidence}
If query-driven RL is learning anything useful, it should beat pure randomization. \textbf{It does, by a small but consistent margin:} 5.444 vs.\ 5.406 on average, and RL sits above the randomized trajectory in essentially every one of the 36 weeks (figure, left). The gap is modest ($\sim$0.04 on the raw CAE scale, $\sim$1\% relative) but persistent rather than noise-driven.
\end{alertblock}

\note{
Be explicit that these are simulation results on an MRT-calibrated generative model, not trial results.

This is the reward-shaping agent at $\bar\gamma{=}0.5$ (\texttt{rs\_g05}) vs.\ the randomized $p=0.5$ baseline (\texttt{random\_send}) --- no averaging across other RL variants or design axes.

Numbers: 100 seeds $\times$ 100 users on \texttt{env\_para\_positivenegative}, aggregated 2026-07-07. \% sent is the mean action propensity, not the realized action count. $^*$ Belief-state $\widehat e_w$ trajectories are only logged for the randomized reference arm in the current (compact-mode) runs; getting it for the RL arm requires a full-mode rerun.
}
\end{frame}
